\documentclass[a4paper,11pt]{article}
\usepackage[utf8]{inputenc}
\usepackage[french]{babel}
\usepackage[T1]{fontenc}
\usepackage{amsmath,amssymb,amsthm}
\usepackage{geometry}
\geometry{margin=2cm}
\usepackage{enumitem}
\usepackage{hyperref}
\usepackage{booktabs}
\usepackage{array}
\usepackage{xcolor}
\usepackage{listings}
\usepackage{titlesec}
\usepackage{fancyhdr}
\usepackage{lastpage}
\usepackage{graphicx}
\usepackage{etoolbox}
\AtBeginDocument{\shorthandoff{;:!?}}

\titleformat{\section}{\large\bfseries\color{blue!60!black}}{}{0pt}{}[\vspace{-0.5ex}\rule{\textwidth}{0.4pt}]
\titleformat{\subsection}{\bfseries\color{blue!40!black}}{}{0pt}{}
\titleformat{\subsubsection}{\itshape\color{blue!30!black}}{}{0pt}{}

\lstdefinestyle{monstyle}{
  frame=single,
  numbers=left,
  numbersep=5pt,
  breaklines=true,
  basicstyle=\small\ttfamily,
  backgroundcolor=\color{gray!5},
  rulecolor=\color{gray!40},
  columns=fullflexible,
  keepspaces=true,
  showstringspaces=false
}
\lstset{style=monstyle}

\pagestyle{fancy}
\fancyhf{}
\fancyhead[L]{\small STA211 -- M\'ethodes supervis\'ees (fondamentaux)}
\fancyhead[R]{\small Fiche r\'ecapitulative}
\fancyfoot[C]{\thepage/\pageref{LastPage}}
\renewcommand{\headrulewidth}{0.4pt}

\theoremstyle{definition}
\newtheorem*{definition}{D\'efinition}
\theoremstyle{remark}
\newtheorem*{remark}{Remarque}
\newtheorem*{important}{\`A retenir}

\title{\textbf{Fiche r\'ecapitulative -- STA211}\\[4pt]
\large \'El\'ements fondamentaux pour m\'ethodes supervis\'ees \\[2pt]
\small Apprentissage statistique, compromis biais-variance,
choix de mod\`eles et validation}
\author{Vincent Audigier, Nd\`eye Niang (CNAM)}
\date{14 juin 2026}

\begin{document}
\maketitle
\thispagestyle{empty}

\begin{abstract}
Cette fiche pr\'esente les fondements des m\'ethodes supervis\'ees
en data-mining. Apr\`es avoir pos\'e le cadre de l'apprentissage
statistique (r\'egression, classification, fonctions de perte), elle
d\'eveloppe le compromis biais-variance, puis les strat\'egies de
choix de mod\`eles : crit\`eres p\'enalis\'es (AIC, BIC), m\'ethodes
pas-\`a-pas, d\'ecoupage apprentissage/test, validation crois\'ee,
et courbe ROC/AUC. Des exemples en R/Python illustrent chaque
concept.
\end{abstract}

\vfill
\tableofcontents
\newpage

% ===================================================================
\section{Introduction}
% ===================================================================

Le \textbf{data-mining} vise \`a identifier des structures au sein
de donn\'ees. On distingue~:
\begin{itemize}
  \item Les \textbf{mod\`eles}~: expliquent et/ou pr\'edisent une
    variable r\'eponse $Y$ \`a partir de variables explicatives $X$.
  \item Les \textbf{patterns}~: r\'egularit\'es locales sans
    variable cible.
\end{itemize}

\subsection{R\'egression \emph{vs} classification}

\begin{itemize}
  \item \textbf{R\'egression}~: $Y$ est continue.
  \item \textbf{Classification supervis\'ee (discrimination)}~:
    $Y$ est qualitative.
\end{itemize}

\subsection{M\'ethodes param\'etriques \emph{vs}
  non-param\'etriques}

\begin{itemize}
  \item \textbf{Param\'etriques}~: forme a priori pour la fonction
    de lien $f$, modul\'ee par des param\`etres
    (ex.~r\'egression lin\'eaire, logistique).
  \item \textbf{Non-param\'etriques}~: aucune forme sp\'ecifique
    suppos\'ee (ex.~arbres, for\^ets al\'eatoires, SVM, r\'eseaux
    de neurones).
\end{itemize}

\begin{table}[h]
\centering
\caption{Typologie des m\'ethodes supervis\'ees}
\begin{tabular}{lcc}
\toprule
& R\'egression & Classification \\
\midrule
Param\'etrique
  & R\'eg.~lin\'eaire
  & R\'eg.~logistique, ALD, mod\`eles additifs g\'en\'eralis\'es \\
Non-param\'etrique
  & Arbres, for\^ets al\'eatoires, r\'eseaux de neurones, splines
  & Arbres, for\^ets al\'eatoires, SVM, r\'eseaux de neurones \\
\bottomrule
\end{tabular}
\end{table}

% ===================================================================
\section{Apprentissage statistique}
% ===================================================================

\subsection{Notations et contexte}

Soit $Y$ la variable r\'eponse et
$\mathbf{X} = (X_1,\dots,X_p)$ le vecteur des $p$ variables
explicatives observ\'e sur $n$ individus~:
\[
\{(\mathbf{x}_1, y_1), \dots, (\mathbf{x}_n, y_n)\}
\in \mathcal{X} \times \mathcal{Y}
\]

\noindent Deux objectifs~:
\begin{itemize}
  \item \textbf{Pr\'ediction}~: pour $\mathbf{x}_0$,
    pr\'edire $\hat{y}_0$ le plus proche possible de $y_0$.
  \item \textbf{Estimation}~: approcher la fonction $f$ qui relie
    $Y$ \`a $X$.
\end{itemize}

\subsection{Fonctions de perte}

\paragraph{R\'egression~:} On suppose $Y = f(X) + \varepsilon$
avec $\mathbb{E}[\varepsilon]=0$ et $\mathbb{V}ar[\varepsilon]=\sigma^2$.

Fonction de contraste des \textbf{moindres carr\'es}~:
\[
\gamma(\hat{f}, (x, y)) = (\hat{f}(x) - y)^2
\]
Perte (esp\'erance)~:
\[
P_\gamma(\hat{f}) = \mathbb{E}_{X,Y}\big[ (Y - \hat{f}(X))^2 \big]
\]

\paragraph{Classification~:} On estime la probabilit\'e
$a~posteriori$ $\mathbb{P}(Y=m_q|X)$.

Contraste \textbf{0-1}~:
\[
\gamma(\hat{f}, (x, y)) = \mathbbm{1}_{\hat{f}(x) \neq y}
\]
Perte (taux d'erreur)~:
\[
P_\gamma(\hat{f}) = \mathbb{P}(\hat{f}(X) \neq Y)
\]

\noindent En estimation, on minimise~:
\[
\mathbb{E}\Big[ \sum_{q=1}^k |\hat{p}(Y=m_q|X) - \mathbb{P}(Y=m_q|X)| \Big]
\]

\subsection{Sur-ajustement (overfitting)}

Un mod\`ele trop complexe (ex.~$k=1$ en k-NN) donne une erreur
d'apprentissage nulle mais ne g\'en\'eralise pas. Il faut un
\'echantillon \textbf{ind\'ependant} pour estimer la perte.

\subsection{Donn\'ees d\'es\'equilibr\'ees}

Cas fr\'equent~: une modalit\'e minoritaire (ex.~30\% d'impay\'es).
Le contraste 0-1 \'etant sym\'etrique, le mod\`ele ignore la
classe rare.

\paragraph{Solutions~:}
\begin{itemize}
  \item \textbf{Contraste asym\'etrique}~:
    $\gamma = w_0 \mathbbm{1}_{\hat{f}(x)\neq0,\,y=0}
      + w_1 \mathbbm{1}_{\hat{f}(x)\neq1,\,y=1}$
  \item \textbf{Sur-\'echantillonnage}~: tirage avec remise dans la
    classe minoritaire.
  \item \textbf{Sous-\'echantillonnage}~: tirage dans la classe
    majoritaire.
  \item \textbf{SMOTE}~: g\'en\`ere des individus fictifs sur des
    segments entre voisins minoritaires (Chawla et al., 2002).
\end{itemize}

\paragraph{Algorithme SMOTE (sur-echantillonnage)~:}
Pour chaque individu $i$ de la classe minoritaire~:
\begin{enumerate}
  \item Identifier ses $k=5$ plus proches voisins (minoritaires).
  \item Tirer un individu $i'$ au hasard parmi les $k$.
  \item Pour $j = 1,\dots,p$~:
    $x_{i''j} \leftarrow x_{ij} + u \times (x_{ij} - x_{i'j})$,
    avec $u \sim \mathcal{U}[0,1]$.
\end{enumerate}

% ===================================================================
\section{Compromis biais - variance}
% ===================================================================

La qualit\'e d'un mod\`ele est li\'ee \`a sa complexit\'e.
Un mod\`ele trop simple (biais \'elev\'e) sous-ajuste, un mod\`ele
trop complexe (variance \'elev\'ee) sur-ajuste.

\subsection{D\'ecomposition de l'erreur}

En r\'egression (MSE), l'erreur esp\'er\'ee se d\'ecompose~:
\[
\mathbb{E}\big[ (Y - \hat{f}(X))^2 \big] =
\underbrace{\mathbb{V}ar[\hat{f}(X)]}_{\text{variance}}
+ \underbrace{(\mathbb{E}[\hat{f}(X)] - f(X))^2}_{\text{biais}^2}
+ \underbrace{\sigma^2}_{\text{bruit}}
\]

\subsection{Illustration (k-NN sur German Credit)}

Avec $k$ petit (mod\`ele complexe)~: erreur train faible,
erreur test \'elev\'ee (surapprentissage).
Avec $k$ optimal ($k\approx10$)~: meilleur compromis.
Au-del\`a~: erreur test augmente (biais trop fort).

\begin{figure}[h]
\centering
\fbox{\parbox{0.7\textwidth}{\centering
  \textbf{Taux de mauvais classement} $\downarrow$\\
  \begin{tabular}{lcc}
  $k$ & Erreur train & Erreur test \\
  \midrule
  1   & 0\%          & $\sim$45\% \\
  10  & $\sim$25\%   & $\sim$30\% \\
  50  & $\sim$35\%   & $\sim$35\% \\
  \end{tabular}
}}
\caption{Compromis biais-variance pour k-NN sur German Credit}
\end{figure}

% ===================================================================
\section{Choix de mod\`eles}
% ===================================================================

\subsection{Crit\`eres p\'enalis\'es}

On ajoute une p\'enalit\'e de complexit\'e au risque empirique~:
\[
\mathcal{L}(\hat{f}) = \frac{1}{n}\sum_{i=1}^n
\gamma(\hat{f}, (\mathbf{x}_i, y_i)) + \text{p\'enalit\'e}
\]

\paragraph{AIC (Akaike Information Criterion)}~:
\[
\text{AIC} = -2\ln L(\hat{\theta}) + 2k
\]
\noindent o\`u $k$ est le nombre de param\`etres.
\'A privil\'egier pour la \textbf{pr\'ediction}.

\paragraph{BIC (Bayesian Information Criterion)}~:
\[
\text{BIC} = -2\ln L(\hat{\theta}) + \ln(n)\,k
\]
P\'enalit\'e plus forte que l'AIC ($\ln(n) > 2$ pour $n \ge 8$).
\'A privil\'egier pour l'\textbf{estimation} (s\'election du vrai
mod\`ele asymptotiquement).

\paragraph{M\'ethodes pas-\`a-pas~:}
\begin{itemize}
  \item \textbf{Descendante}~: part du mod\`ele complet, \'elimine
    les variables une \`a une.
  \item \textbf{Ascendante}~: part du mod\`ele vide, ajoute les
    variables une \`a une.
  \item \textbf{Progressive}~: ascendante avec possibilit\'e
    d'\'elimination.
\end{itemize}

\paragraph{Limites des crit\`eres p\'enalis\'es~:}
\begin{itemize}
  \item N\'ecessitent un mod\`ele param\'etrique (vraisemblance).
  \item Ne comparent pas des mod\`eles de familles diff\'erentes.
  \item Le nombre de param\`etres ne refl\`ete pas toujours la
    complexit\'e (ex.~r\'egression ridge).
\end{itemize}

\subsection{Approches empiriques}

\paragraph{D\'ecoupage test/apprentissage~:}
\begin{enumerate}
  \item S\'eparer les donn\'ees en $\mathcal{A}$ (apprentissage)
    et $\mathcal{T}$ (test).
  \item Ajuster chaque mod\`ele sur $\mathcal{A}$.
  \item Choisir celui qui minimise l'erreur sur $\mathcal{T}$.
  \item R\'e-ajuster le mod\`ele retenu sur
    $\mathcal{A} \cup \mathcal{T}$.
\end{enumerate}

\paragraph{Validation crois\'ee $K$-fold~:}
\begin{enumerate}
  \item D\'ecouper l'\'echantillon en $K$ blocs de m\^eme taille.
  \item Pour chaque bloc~: les $K-1$ autres servent
    d'apprentissage, le bloc sert de test.
  \item Moyenner les $K$ erreurs de test.
\end{enumerate}
$K=10$ est le choix usuel. $K=n$ = \textbf{leave-one-out} (LOO).

\paragraph{Mesure objective de la perte~:}
3 \'echantillons~: \textbf{apprentissage} (ajuster),
\textbf{validation} (choisir), \textbf{test} (mesure finale).
Proportion typique~: $1/2$, $1/4$, $1/4$.

\subsection{Compl\'ement pour la classification~: ROC et AUC}

Pour une variable binaire, on peut faire varier le seuil $s$ de
d\'ecision $\hat{p}(Y=1|X) > s$.

\textbf{Courbe ROC}~: taux de vrais positifs en fonction du taux
de faux positifs pour tous les seuils $s \in [0,1]$.

\textbf{AUC (Area Under the Curve)}~:
\begin{itemize}
  \item AUC $=0.5$~: mod\`ele al\'eatoire (aucun pouvoir pr\'edicif).
  \item AUC $=1$~: mod\`ele parfait.
  \item En pratique~: AUC $>0.7$ acceptable, $>0.8$ excellent.
\end{itemize}

% ===================================================================
\section{Code Python}
% ===================================================================

\subsection{Validation crois\'ee avec scikit-learn}

\begin{lstlisting}[language=Python]
from sklearn.model_selection import cross_val_score, train_test_split
from sklearn.neighbors import KNeighborsClassifier
from sklearn.datasets import make_classification
import numpy as np

X, y = make_classification(n_samples=200, n_features=2,
                           n_redundant=0, random_state=42)
X_train, X_test, y_train, y_test = train_test_split(
    X, y, test_size=0.33, random_state=42)

for k in [1, 5, 10, 15, 30]:
    model = KNeighborsClassifier(n_neighbors=k)
    scores = cross_val_score(model, X_train, y_train, cv=10)
    model.fit(X_train, y_train)
    test_score = model.score(X_test, y_test)
    print(f"k={k:2d} | CV: {scores.mean():.3f} (+/-{scores.std():.2f}) "
          f"| Test: {test_score:.3f}")
\end{lstlisting}

\subsection{SMOTE avec imbalanced-learn}

\begin{lstlisting}[language=Python]
from imblearn.over_sampling import SMOTE
from collections import Counter

smote = SMOTE(random_state=42)
X_res, y_res = smote.fit_resample(X_train, y_train)
print("Avant SMOTE:", Counter(y_train))
print("Apres SMOTE:", Counter(y_res))
\end{lstlisting}

\subsection{Courbe ROC et AUC}

\begin{lstlisting}[language=Python]
from sklearn.metrics import roc_curve, auc
import matplotlib.pyplot as plt

model = KNeighborsClassifier(n_neighbors=10)
model.fit(X_train, y_train)
y_prob = model.predict_proba(X_test)[:, 1]
fpr, tpr, _ = roc_curve(y_test, y_prob)
roc_auc = auc(fpr, tpr)

plt.plot(fpr, tpr, label=f"AUC = {roc_auc:.3f}")
plt.plot([0,1], [0,1], 'k--')
plt.xlabel("Taux faux positifs")
plt.ylabel("Taux vrais positifs")
plt.legend()
plt.show()
\end{lstlisting}

% ===================================================================
\section{Code R}
% ===================================================================

\subsection{Validation crois\'ee avec caret}

\begin{lstlisting}[language=R]
library(caret)
library(class)

data(GermanCredit)
train_idx <- createDataPartition(GermanCredit$Class,
                                 p = 0.67, list = FALSE)
train <- GermanCredit[train_idx, ]
test  <- GermanCredit[-train_idx, ]

ctrl <- trainControl(method = "cv", number = 10)
for (k in c(1, 5, 10, 15, 30)) {
  model <- train(Class ~ ., data = train, method = "knn",
                 trControl = ctrl, tuneGrid = data.frame(k = k))
  cat(sprintf("k=%2d | CV: %.3f | Test: %.3f\n",
              k, max(model$results$Accuracy),
              predict(model, test) |> (\(x)
                mean(x == test$Class))()))
}
\end{lstlisting}

\subsection{SMOTE avec le package UBL}

\begin{lstlisting}[language=R]
library(UBL)
smote_train <- SmoteClassif(Class ~ ., train, C.perc = "balance")
table(smote_train$Class)
\end{lstlisting}

\subsection{Courbe ROC}

\begin{lstlisting}[language=R]
library(pROC)
model <- knn3(Class ~ ., train, k = 10)
prob  <- predict(model, test)[, 2]
roc_obj <- roc(test$Class, prob)
plot(roc_obj, print.auc = TRUE)
\end{lstlisting}

% ===================================================================
\newpage
\section*{Questions cl\'es (QCM)}
\addcontentsline{toc}{section}{Questions cl\'es (QCM)}

\begin{enumerate}

  \item Quelle est la fonction de contraste usuelle en r\'egression~?
    \begin{enumerate}
      \item Contraste 0-1
      \item Contraste des moindres carr\'es
      \item Contraste log-vraisemblance
      \item Contraste asym\'etrique
    \end{enumerate}
    \textbf{R\'eponse~: b (moindres carr\'es).}

  \item En classification, la perte $P_\gamma(\hat{f}) =
    \mathbb{P}(\hat{f}(X) \neq Y)$ correspond \`a~:
    \begin{enumerate}
      \item L'AUC
      \item Le taux de mauvais classement
      \item La somme des carr\'es r\'esiduels
      \item Le BIC
    \end{enumerate}
    \textbf{R\'eponse~: b (taux d'erreur).}

  \item Que se passe-t-il quand $k=1$ en k-NN sur l'erreur
    d'apprentissage~?
    \begin{enumerate}
      \item Erreur maximale
      \item Erreur nulle (sur-apprentissage)
      \item Erreur \'egale \`a l'erreur de test
      \item Le mod\`ele ne converge pas
    \end{enumerate}
    \textbf{R\'eponse~: b (erreur nulle, chaque point est son propre
    plus proche voisin).}

  \item La proc\'edure SMOTE permet de~:
    \begin{enumerate}
      \item R\'eduire la dimension
      \item G\'en\'erer des individus fictifs pour la classe
        minoritaire
      \item S\'electionner des variables
      \item Calculer l'AUC
    \end{enumerate}
    \textbf{R\'eponse~: b (sur-\'echantillonnage synth\'etique).}

  \item Dans la d\'ecomposition biais-variance, un mod\`ele trop
    simple a~:
    \begin{enumerate}
      \item Un biais \'elev\'e, une variance faible
      \item Un biais faible, une variance \'elev\'ee
      \item Un biais \'elev\'e, une variance \'elev\'ee
      \item Un biais nul, une variance nulle
    \end{enumerate}
    \textbf{R\'eponse~: a (sous-apprentissage).}

  \item Le crit\`ere AIC p\'enalise la log-vraisemblance par~:
    \begin{enumerate}
      \item $2k$
      \item $\ln(n)\,k$
      \item $k$
      \item $\sqrt{k}$
    \end{enumerate}
    \textbf{R\'eponse~: a (AIC = $-2\ln L + 2k$).}

  \item Pour $n=100$, quel crit\`ere p\'enalise le plus la
    complexit\'e~?
    \begin{enumerate}
      \item AIC
      \item BIC
      \item Les deux p\'enalisent autant
      \item Aucun des deux
    \end{enumerate}
    \textbf{R\'eponse~: b (BIC~: $\ln(100) \approx 4.6 > 2$).}

  \item En validation crois\'ee $K$-fold, que vaut $K$ pour le
    leave-one-out~?
    \begin{enumerate}
      \item $K=10$
      \item $K=5$
      \item $K=n$
      \item $K=1$
    \end{enumerate}
    \textbf{R\'eponse~: c ($K=n$, chaque individu est son propre
    bloc de test).}

  \item Une AUC de 0.5 signifie que le mod\`ele~:
    \begin{enumerate}
      \item Est parfait
      \item Est al\'eatoire (aucun pouvoir pr\'edictif)
      \item Est excellent
      \item A un biais nul
    \end{enumerate}
    \textbf{R\'eponse~: b (AUC=0.5 = pr\'ediction al\'eatoire).}

  \item Quel est l'ordre des 3 \'echantillons pour une mesure
    objective de la perte~?
    \begin{enumerate}
      \item Test $\to$ Validation $\to$ Apprentissage
      \item Apprentissage $\to$ Validation $\to$ Test
      \item Apprentissage $\to$ Test $\to$ Validation
      \item Validation $\to$ Apprentissage $\to$ Test
    \end{enumerate}
    \textbf{R\'eponse~: b (apprentissage~: ajuster, validation~:
    choisir, test~: mesurer).}

\end{enumerate}

% ===================================================================
\begin{thebibliography}{9}
% ===================================================================

\bibitem{Chawla2002}
  N.~V.~Chawla, K.~W.~Bowyer, L.~O.~Hall, W.~P.~Kegelmeyer.
  \newblock \emph{SMOTE: Synthetic Minority Over-sampling Technique}.
  \newblock J.~Artif.~Intell.~Res., 16:321--357, 2002.

\bibitem{Hastie2009}
  T.~Hastie, R.~Tibshirani, J.~Friedman.
  \newblock \emph{The Elements of Statistical Learning}.
  \newblock Springer, 2009.

\bibitem{Saporta2006}
  G.~Saporta.
  \newblock \emph{Probabilit\'es, analyse des donn\'ees et
    statistique}.
  \newblock Editions Technip, 2006.

\bibitem{Cornillon2010}
  P.-A.~Cornillon, E.~Matzner-Lober.
  \newblock \emph{R\'egression avec R}.
  \newblock Springer, 2010.

\bibitem{Tuffery2007}
  S.~Tuff\'ery.
  \newblock \emph{Data Mining et Statistique d\'ecisionnelle}.
  \newblock Editions Technip, 2007.

\end{thebibliography}

\end{document}
